% ============================================================
% EJERCICIO 8 — Artículo académico avanzado (Módulo 2)
% 4-5 páginas, IA en TODAS las etapas (título, resumen, ecuaciones, tablas con
% datos SIMULADOS declarados como tales, diagrama TikZ, corrección de estilo,
% formato de bibliografía). Temas sugeridos: convertidores DC-DC, filtro activo
% biomédico, MPPT solar, PID embebido.
% Para el formato de referencias use APA 7: ver apa7/articulo_apa7.tex
% ============================================================
\documentclass[11pt,a4paper,twocolumn]{article}
\usepackage[utf8]{inputenc}
\usepackage[T1]{fontenc}
\usepackage[spanish]{babel}
\usepackage{amsmath,amssymb}
\usepackage{booktabs}
\usepackage{siunitx}
\sisetup{locale=FR, output-decimal-marker={,}}
\usepackage{circuitikz}
\usepackage[hidelinks]{hyperref}
\usepackage[margin=2cm]{geometry}

\title{TODO Título}
\author{TODO Autor\\ Universidad Nacional Tecnológica de Lima Sur}
\date{\today}

\begin{document}
\twocolumn[
  \maketitle
  \begin{abstract}
    % TODO 150-200 palabras: problema, método, resultados, conclusión.
  \end{abstract}
  \textbf{Palabras clave:} TODO, TODO, TODO
  \vspace{1em}
]

\section{Introducción}
\section{Marco teórico}
\section{Metodología}
% TODO ecuaciones + diagrama (circuitikz/tikz)
\section{Resultados}
% TODO tabla con datos simulados (declare que son simulados) y gráfica
\section{Conclusiones}

% TODO bibliografía (APA 7 con biblatex, ver apa7/) con fuentes verificadas

\appendix
\section{Declaración de uso de IA}
% TODO herramienta, etapa, prompt, verificación realizada
\end{document}
% RÚBRICA (20 pts): estructura y redacción (5) · ecuaciones/figuras/tablas (5)
% · bibliografía verificada y bien formateada (4) · uso de IA documentado (4)
% · compila sin avisos (2)
